\subsection{Orthogonal bases.}

DEFINITION 1. --- Let Φ be a Hermitian form on E. A basis ( \(e_{i}\) ) of E is said to be orthogonal for Φ if any two elements of this basis are orthogonal for Φ.

If moreover \(\Phi(e_i, e_i) = 1\) for all \(i\) , the basis \((e_i)\) is said to be orthonormal.

Let \((e_i)\) be an orthogonal basis; if we set \(\Phi(e_i, e_i) = \alpha_i\) , one has

\[
\Phi (\sum_ {i} \xi_ {i} e _ {i}, \sum_ {i} \eta_ {i} e _ {i}) = \sum_ {i} \xi_ {i} \alpha_ {i} \overline {{\eta}} _ {i}.\tag{1}
\]

Lemma 1. --- Suppose that A is a field and \(\Phi\) a Hermitian form \(\neq 0\) on E. If all vectors of E are isotropic, then A is a commutative field of characteristic 2, the antiautomorphism J is the identity, and \(\Phi\) is alternating.

Indeed, if one expands \(\Phi(x+y, x+y)=0\) , one obtains, taking the hypotheses into account, \(\Phi(x,x)=\Phi(y,y)=0\) , the relation \(\Phi(x,y)=-\overline{\Phi(x,y)}\) for all \(x,y\) in E. Since \(\Phi\) is \(\neq 0\) , there exist \(x,y\) in E such that \(\Phi(x,y)=1\) . Writing that \(\Phi(\lambda x,y)=-\overline{\Phi(\lambda x,y)}\) , one obtains \(\bar{\lambda}=-\lambda\) for all \(\lambda\in A\) . Taking first \(\lambda=1\) , one sees that A is of characteristic 2; the relation \(\bar{\lambda}=-\lambda\) then shows that J is the identity, hence A is commutative and \(\Phi\) is alternating.

THEOREM 1. --- Suppose that A is a field and E a vector space of finite dimension n over A. Then, for every Hermitian form Φ on E, E admits an orthogonal basis, except if the following conditions are simultaneously satisfied:

\begin{enumerate}
\def\labelenumi{(\Alph{enumi})}
\setcounter{enumi}{2}
\tightlist
\item
  A is commutative of characteristic 2, the antiautomorphism J is the identity, Φ is alternating and nonzero.
\end{enumerate}

Reason by induction on \(n\) , the result being evident for \(n=0\) . We may suppose \(\Phi\neq0\) . If (C) is not satisfied, Lemma 1 shows that there exists an element \(x\in\mathbf{E}\) such that \(\Phi(x,x)\neq0\) . Let H be the subspace of E orthogonal to \(x\) ; it has dimension \(\geqslant n-1\) , and, since \(x \notin H\) , \(H\) is exactly of dimension \(n-1\) . If the restriction \(\Psi\) of \(\Phi\) to \(H\) does not satisfy (C), then by the induction hypothesis there exists a basis \((e_2, \ldots, e_n)\) of \(H\) which is orthogonal for \(\Psi\) ; by setting \(e_1 = x\) one obtains an orthogonal basis \((e_1, e_2, \ldots, e_n)\) of \(E\) . It remains to examine the case where \(A\) is a commutative field of characteristic 2, where \(J\) is the identity, and where \(\Psi\) is alternating and \(\neq 0\) . There then exists \(y\) , \(z\) in \(H\) such that \(\Psi(y, z) \neq 0\) ; set \(e_1 = x + y\) ; for \(x + \lambda z (\lambda \in A)\) to be orthogonal to \(e_1\) , it is necessary and sufficient that we have \(0 = \Phi(x + y, x + \lambda z) = \Phi(x, x) + \lambda \Psi(y, z)\) , a condition which determines \(\lambda\) in a unique way; the scalar \(\lambda\) having been chosen thus, we have \(\Phi(x + \lambda z, x + \lambda z) = \Phi(x, x) \neq 0\) , hence the restriction \(\Psi'\) of \(\Phi\) to the subspace \(H'\) of \(E\) orthogonal to \(e_1\) is not alternating; one may consequently apply the induction hypothesis to \(H'\) , which proves the theorem.

When (C) is satisfied, there evidently exists no orthogonal basis for \(\Phi\) .

COROLLARY 1. --- With the notations of th. 1, one supposes moreover that (C) is not satisfied, that \(\Phi\) is nondegenerate and that, for every \(x \in E\) , there exists \(\rho \in A\) such that \(\Phi(x, x) = \rho \bar{\rho}\) . Then E admits an orthonormal basis for \(\Phi\) .

Indeed, let \((e_i)\) ( \(i = 1, \ldots, n\) ) be an orthogonal basis of E. Put \(\Phi(e_i, e_i) = \alpha_i\) . We have \(\alpha_i \neq 0\) for \(i = 1, \ldots, n\) since \(\Phi\) is nondegenerate. By hypothesis, there exist elements \(\beta_i\) of A such that \(\alpha_i = \beta_i \overline{\beta_i}\) for \(i = 1, \ldots, n\) ; we have \(\beta_i \neq 0\) . By setting \(f_i = \beta_i^{-1} e_i\) , one has \(\Phi(f_i, f_i) = \beta_i^{-1} \alpha_i \overline{\beta_i^{-1}} = 1\) for all \(i\) , and \(\Phi(f_i, f_j) = 0\) for \(i \neq j\) . Hence \((f_i)\) is an orthonormal basis.

Remark. --- The last hypothesis of the corollary is satisfied when J is the identity and every element of A is the square of an element of A (for example, when A is algebraically closed).

COROLLARY 2. --- Let A be a field and R a Hermitian matrix of order n and rank r over A. Then, unless the following condition holds:

(C') A is commutative of characteristic 2, J is the identity, R is alternating and nonzero,

there exists an invertible matrix P of order n over A such that

\[
{ } ^ { t } P . R . \overline { { P } } = \left( \begin{array} { c c c c c } \alpha _ { 1 } & 0 \dots 0 \dots 0 \\ 0 & \alpha _ { 2 } \dots 0 \dots 0 \\ \cdot & \cdot & \cdot & \cdot & \cdot \\ 0 & 0 \dots \alpha _ { r } \dots 0 \\ 0 & 0 \dots 0 \dots 0 \\ \cdot & \cdot & \cdot & \cdot & \cdot \\ 0 & 0 \dots 0 \dots 0 \end{array} \right)
\]

\(where \overline{\alpha}_{i} = \alpha_{i} \neq 0 \text{ for } i = 1, \ldots, r.\)

PROPOSITION 1. --- Suppose that A is a commutative field. Let \(\Phi\) be a Hermitian form on E, and \((x_{n})\) ( \(n=1,2,\ldots\) ) a sequence (finite or infinite) of linearly independent vectors of E such that, for every \(n\) , the subspace \(\mathrm{E}_{n}=\mathrm{A}x_{1}+\cdots+\mathrm{A}x_{n}\) is non-isotropic. Let \(\mathrm{D}_{jn}\) ( \(j\leqslant n\) ) be the cofactor of \(\Phi(x_{j},x_{n})\) in the matrix \((\Phi(x_{s},x_{t}))_{(s,t=1,\ldots,n)}\) . We then have \(\mathrm{D}_{nn}\neq0\) for all \(n\) . Set

\[
e _ {n} = \sum_ {j = 1} ^ {n} \mathrm{D} _ {n n} ^ {- 1} \mathrm{D} _ {j n} x _ {j}.\tag{2}
\]

Then, for every \(n\) , \((e_{1},\ldots,e_{n})\) is an orthogonal basis of \(\mathrm{E}_{n}\) and we have

\[
\Phi (e _ {n}, e _ {n}) = \mathrm{D} _ {n n} ^ {- 1} \mathrm{D} _ {n + 1, n + 1}.\tag{3}
\]

Indeed, since the restriction of \(\Phi\) to \(\mathbf{E}_{n-1}\) is nondegenerate, we have \(\mathrm{D}_{nn} \neq 0\) ( \(\S 2\) , prop. 3); note that we have \(\mathrm{D}_{11} = 1\) since the determinant of the empty matrix is equal to 1. Formulas (2) imply first that we have \(e_n \equiv x_n\) (mod. \(\mathbf{E}_{n-1}\) ) for all \(n\) , hence that the \(e_n\) are linearly independent, and that \((e_1, \ldots, e_n)\) is a basis of \(\mathbf{E}_n\) . For every \(j < n\) , one has

\[
\Phi (e _ {n}, x _ {j}) = \mathrm{D} _ {n n} ^ {- 1} \sum_ {k = 1} ^ {n} \mathrm{D} _ {k n} \Phi (x _ {k}, x _ {j}) = 0
\]

(chap. III, § 6, no. 1, formula (12)); hence \(e_n\) is orthogonal to \(\mathrm{E}_{n-1}\) , and in particular to \(e_j\) for \(j < n\) . On the other hand, we have

\[
\begin{array}{r l} \Phi (e _ {n}, e _ {n}) & = \Phi (e _ {n}, \sum_ {j = 1} ^ {n} \mathrm{D} _ {n n} ^ {- 1} \mathrm{D} _ {j n} x _ {j}) = \Phi (e _ {n}, x _ {n}) = \Phi (\sum_ {j = 1} ^ {n} \mathrm{D} _ {n n} ^ {- 1} \mathrm{D} _ {j n} x _ {j}, x _ {n}) \\ & = \mathrm{D} _ {n n} ^ {- 1} \sum_ {j = 1} ^ {n} \mathrm{D} _ {j n} \Phi (x _ {j}, x _ {n}) = \mathrm{D} _ {n n} ^ {- 1} \mathrm{D} _ {n + 1, n + 1} \end{array}
\]

(chap. III, § 6, no. 1, formula (10)). This proves our assertions.

With the notations of prop. 1, one says that the sequence \((e_n)\) is obtained from the sequence \((x_n)\) by the Gram-Schmidt orthogonalization process.

PROPOSITION 2. --- Let \(\Phi\) be a Hermitian form on E, and \((e_i)\) ( \(i = 1, \ldots, n\) ) an orthogonal basis (resp. orthonormal basis) of E for \(\Phi\) . Then, for every \(p \geqslant 0\) , the basis of \(\bigotimes^{p} \mathbf{E}\) formed by the \(e_{i_1} \otimes \cdots \otimes e_{i_p}\) and the basis \((e_h)\) of \(\bigwedge^{p} \mathbf{E}\) (where H ranges over the set of subsets with \(p\) elements of {[}1, n{]}; cf. chap. III, § 5, no. 6) are orthogonal (resp. orthonormal) for the extensions of \(\Phi\) to \(\bigotimes^{p} \mathbf{E}\) and \(\bigwedge^{p} \mathbf{E}\) respectively (§ 1, no. 9). If, moreover, the maps associated with \(\Phi\) are bijective, the basis \((e_i')\) of \(\mathbf{E}^*\) dual of \((e_i)\) is orthogonal (resp. orthonormal) for the inverse form \(\hat{\Phi}\) of \(\Phi\) (§ 1, no. 7).

The assertions concerning \(\widehat{\otimes}\) E and \(\wedge\) E follow immediately from formulas (35) and (37) of § 1, no. 9. The one relating to the inverse form follows from the fact that the matrix of \(\widehat{\Phi}\) with respect to \((e_i')\) is the inverse of the matrix of \(\Phi\) with respect to \((e_i)\) ( \(§ 1\) , no. 10).

